\chapter{Marco General de Ciberseguridad NIST CSF 2.0}
\label{ch:nist_csf}

En este capítulo se articula el gobierno de la ciberseguridad universitaria bajo el estándar internacional \textbf{NIST Cybersecurity Framework 2.0}, integrando sus seis funciones fundamentales.

\section{Las Seis Funciones Fundamentales del NIST CSF 2.0 en la UNAP}
\begin{enumerate}
    \item \textbf{GOVERN (Gobernar):} Establecimiento y difusión de directivas institucionales de seguridad digital, designación formal del Oficial de Seguridad de la Información (CISO), asignación presupuestal específica y alineamiento con la normativa del Estado Peruano (D.U. N° 006-2020).
    \item \textbf{IDENTIFY (Identificar):} Catalogación y actualización continua del inventario de activos, evaluación sistemática de riesgos informáticos y análisis de impacto en la operación académica universitaria.
    \item \textbf{PROTECT (Proteger):} Hardening de servidores Nginx, publicación forzada de cabeceras HSTS, configuración de registros SPF estricto (\texttt{-all}) y DMARC (\texttt{p=reject}), y adopción de autenticación multifactor (MFA).
    \item \textbf{DETECT (Detectar):} Monitoreo operacional 24/7 mediante la plataforma SIEM centralizada (Wazuh), correlación de anomalías, alertas automáticas por fuerza bruta y recolección inmutable de registros de eventos.
    \item \textbf{RESPOND (Responder):} Protocolos y playbooks operativos para contención inmediata, aislamiento automático de subredes comprometidas y comunicación institucional con el CSIRT Nacional de la PCM.
    \item \textbf{RECOVER (Recuperar):} Políticas de respaldo bajo el esquema 3-2-1 con copias de seguridad inmutables fuera de línea (\textit{Air-Gapped}) y ensayos periódicos semestrales de restauración de bases de datos.
\end{enumerate}

\section{Matriz NIST CSF 2.0 de la UNAP}

\begin{center}
\small
\begin{longtable}{>{\bfseries}p{2.5cm} p{4.2cm} p{3.8cm} p{4.5cm}}
\toprule
\rowcolor{outermarginbgcolor}
Función NIST & Acción Estratégica & Responsable & Herramientas / Controles \\
\midrule
\endfirsthead
\toprule
\rowcolor{outermarginbgcolor}
Función NIST & Acción Estratégica & Responsable & Herramientas / Controles \\
\midrule
\endhead
\bottomrule
\endfoot
GOVERN & Políticas y directivas de seguridad & Rectorado / CISO & Comité de Gobierno Digital y normativas \\
\midrule
IDENTIFY & Inventario continuo y matriz de riesgos & Unidad de Seguridad & crt.sh, Nessus y matriz de activos \\
\midrule
PROTECT & Hardening web, MFA y DMARC & OTI / Redes & Nginx, Let's Encrypt, TOTP, DNS BIND \\
\midrule
DETECT & Monitoreo continuo de eventos y logs & CSIRT-UNAP & Wazuh SIEM, Suricata IDS y Syslog \\
\midrule
RESPOND & Playbooks de contención de incidentes & CSIRT-UNAP & Procedimientos formales y scripts de aislamiento \\
\midrule
RECOVER & Restauración de respaldos 3-2-1 & OTI / Soporte & Almacenamiento Air-Gapped y Veeam Backup \\
\bottomrule
\end{longtable}
\end{center}